
\documentclass{article}
\usepackage{colortbl}
\usepackage{makecell}
\usepackage{multirow}
\usepackage{supertabular}

\begin{document}

\newcounter{utterance}

\centering \large Interaction Transcript for game `dond', experiment `coop\_en', episode 19 with reflexion{-}Qwen3.5{-}9B.
\vspace{24pt}

{ \footnotesize  \setcounter{utterance}{1}
\setlength{\tabcolsep}{0pt}
\begin{supertabular}{c@{$\;$}|p{.15\linewidth}@{}p{.15\linewidth}p{.15\linewidth}p{.15\linewidth}p{.15\linewidth}p{.15\linewidth}}
    \# & $\;$A & \multicolumn{4}{c}{Game Master} & $\;\:$B\\
    \hline

    \theutterance \stepcounter{utterance}  
    & & \multicolumn{4}{p{0.6\linewidth}}{
        \cellcolor[rgb]{0.9,0.9,0.9}{
            \makecell[{{p{\linewidth}}}]{
                \texttt{\tiny{[P1$\langle$GM]}}
                \texttt{You are playing a cooperative negotiation game in which you have to agree on how to divide a set of items among you and another player.} \\
\\ 
\texttt{Rules:} \\
\texttt{(a) You and the other player are given a set of items. Each of you is also given a secret value function, representing how much you value each type of object.} \\
\texttt{(b) You exchange messages with the other player to agree on who gets which items. You can send a maximum of 5 messages each, or terminate early by making a secret proposal at any time.} \\
\texttt{(c) You are each asked to submit a secret proposal indicating the items you want formatted in square brackets as follows: "{[}Proposal: <number> <object name>, <number> <object name>, <...>{]}"} \\
\texttt{(d) If your proposals are complementary, i.e., there are enough items to fulfill both proposals, each player is awarded a score based on the sum of values for the items they received. Otherwise, both of you get zero points.} \\
\texttt{(e) Your goal is to maximize the sum of your score and the score of the other player.} \\
\\ 
\texttt{Let us start.} \\
\\ 
\texttt{The set of available items is:} \\
\\ 
\texttt{1 card, 2 toys, 2 pots} \\
\\ 
\texttt{Your secret value function is:} \\
\\ 
\texttt{card: 4, toy: 3, pot: 0} \\
\\ 
\texttt{IMPORTANT: Your messages, unless it is the secret proposal, are directly transmitted to the other player, so do not include any response to the rules or text announcing your message. To make a secret proposal, use the indicated format. Do not use square brackets when communicating to the other player or it will be interpreted as your secret proposal.} \\
            }
        }
    }
    & \\ \\

    \theutterance \stepcounter{utterance}  
    & \multicolumn{4}{p{0.6\linewidth}}{
        \cellcolor[rgb]{0.9,0.9,0.9}{
            \makecell[{{p{\linewidth}}}]{
                \texttt{\tiny{[P1$\rangle$GM]}}
                \texttt{I understand the rules and my valuation function. I am ready to negotiate.} \\
\\ 
\texttt{My valuations are:} \\
\texttt{{-} Card: 4} \\
\texttt{{-} Toy: 3} \\
\texttt{{-} Pot: 0} \\
\\ 
\texttt{Since I value pots at 0, I am indifferent to who gets them. I value the card most (4) and toys next (3). To maximize the total score (my score + your score), we should ensure the high{-}value items (card and toys) are allocated to the person who values them most.} \\
\\ 
\texttt{I will start by proposing a division that gives me the card and one toy, while leaving the other items for you, assuming you might value them differently. This is a tentative offer to see your response.} \\
\\ 
\texttt{Here is my first message:} \\
\\ 
\texttt{I propose I take the card and one toy. You can take the remaining toy and both pots. Let me know if this works for you or if you'd like to adjust.} \\
            }
        }
    }
    & & \\ \\

    \theutterance \stepcounter{utterance}  
    & & & \multicolumn{2}{p{0.3\linewidth}}{
        \cellcolor[rgb]{0.9,0.9,0.9}{
            \makecell[{{p{\linewidth}}}]{
                \texttt{\tiny{[GM$|$GM]}}
                \texttt{continue} \\
            }
        }
    }
    & & \\ \\

    \theutterance \stepcounter{utterance}  
    & & & \multicolumn{2}{p{0.3\linewidth}}{
        \cellcolor[rgb]{0.9,0.9,0.9}{
            \makecell[{{p{\linewidth}}}]{
                \texttt{\tiny{[GM$|$GM]}}
                \texttt{I understand the rules and my valuation function. I am ready to negotiate.} \\
\\ 
\texttt{My valuations are:} \\
\texttt{{-} Card: 4} \\
\texttt{{-} Toy: 3} \\
\texttt{{-} Pot: 0} \\
\\ 
\texttt{Since I value pots at 0, I am indifferent to who gets them. I value the card most (4) and toys next (3). To maximize the total score (my score + your score), we should ensure the high{-}value items (card and toys) are allocated to the person who values them most.} \\
\\ 
\texttt{I will start by proposing a division that gives me the card and one toy, while leaving the other items for you, assuming you might value them differently. This is a tentative offer to see your response.} \\
\\ 
\texttt{Here is my first message:} \\
\\ 
\texttt{I propose I take the card and one toy. You can take the remaining toy and both pots. Let me know if this works for you or if you'd like to adjust.} \\
            }
        }
    }
    & & \\ \\

    \theutterance \stepcounter{utterance}  
    & & \multicolumn{4}{p{0.6\linewidth}}{
        \cellcolor[rgb]{0.9,0.9,0.9}{
            \makecell[{{p{\linewidth}}}]{
                \texttt{\tiny{[GM$\rangle$P2]}}
                \texttt{You are playing a cooperative negotiation game in which you have to agree on how to divide a set of items among you and another player.} \\
\\ 
\texttt{Rules:} \\
\texttt{(a) You and the other player are given a set of items. Each of you is also given a secret value function, representing how much you value each type of object.} \\
\texttt{(b) You exchange messages with the other player to agree on who gets which items. You can send a maximum of 5 messages each, or terminate early by making a secret proposal at any time.} \\
\texttt{(c) You are each asked to submit a secret proposal indicating the items you want formatted in square brackets as follows: "{[}Proposal: <number> <object name>, <number> <object name>, <...>{]}"} \\
\texttt{(d) If your proposals are complementary, i.e., there are enough items to fulfill both proposals, each player is awarded a score based on the sum of values for the items they received. Otherwise, both of you get zero points.} \\
\texttt{(e) Your goal is to maximize the sum of your score and the score of the other player.} \\
\\ 
\texttt{Let us start.} \\
\\ 
\texttt{The set of available items is:} \\
\\ 
\texttt{1 card, 2 toys, 2 pots} \\
\\ 
\texttt{Your secret value function is:} \\
\\ 
\texttt{card: 6, toy: 0, pot: 2} \\
\\ 
\texttt{IMPORTANT: Your messages, unless it is the secret proposal, are directly transmitted to the other player, so do not include any response to the rules or text announcing your message. To make a secret proposal, use the indicated format. Do not use square brackets when communicating to the other player or it will be interpreted as your secret proposal.} \\
\\ 
\\ 
\texttt{I understand the rules and my valuation function. I am ready to negotiate.} \\
\\ 
\texttt{My valuations are:} \\
\texttt{{-} Card: 4} \\
\texttt{{-} Toy: 3} \\
\texttt{{-} Pot: 0} \\
\\ 
\texttt{Since I value pots at 0, I am indifferent to who gets them. I value the card most (4) and toys next (3). To maximize the total score (my score + your score), we should ensure the high{-}value items (card and toys) are allocated to the person who values them most.} \\
\\ 
\texttt{I will start by proposing a division that gives me the card and one toy, while leaving the other items for you, assuming you might value them differently. This is a tentative offer to see your response.} \\
\\ 
\texttt{Here is my first message:} \\
\\ 
\texttt{I propose I take the card and one toy. You can take the remaining toy and both pots. Let me know if this works for you or if you'd like to adjust.} \\
            }
        }
    }
    & \\ \\

    \theutterance \stepcounter{utterance}  
    & & & \multicolumn{4}{p{0.6\linewidth}}{
        \cellcolor[rgb]{0.9,0.9,0.9}{
            \makecell[{{p{\linewidth}}}]{
                \texttt{\tiny{[GM$\langle$P2]}}
                \texttt{That sounds like a fair starting point. Since you value the card at 4 and a toy at 3, taking those gives you a score of 7. The remaining items are one toy (which you value at 0) and two pots (which you also value at 0), so your total score would be 0.} \\
\\ 
\texttt{However, looking at my own valuations:} \\
\texttt{{-} Card: 4} \\
\texttt{{-} Toy: 3} \\
\texttt{{-} Pot: 0} \\
\\ 
\texttt{If I take the card and one toy, I get \$4 + 3 = 7\$.} \\
\texttt{If you take the remaining toy and two pots, you get \$0 + 0 = 0\$.} \\
\texttt{Total combined score: \$7 + 0 = 7\$.} \\
\\ 
\texttt{But wait, let's re{-}evaluate the total potential value of the items based on *our* combined valuations to see if we can do better.} \\
\texttt{{-} Card: You value 6, I value 4. Total = 10.} \\
\texttt{{-} Toy: You value 3, I value 0. Total = 3.} \\
\texttt{{-} Pot: You value 2, I value 0. Total = 2.} \\
\\ 
\texttt{The maximum possible total value we can achieve is \$10 + 3 + 2 = 15\$.} \\
\texttt{In your proposal, the total value achieved is only 7. We are leaving a lot of value on the table because the card is going to me (value 4) instead of you (value 6), and the pots are going to you (value 2) instead of me (value 0, so no loss there, but the card is the big miss).} \\
\\ 
\texttt{Actually, let's look at the allocation again.} \\
\texttt{Your proposal:} \\
\texttt{{-} Me: Card, 1 Toy. (My value: \$4 + 3 = 7\$)} \\
\texttt{{-} You: 1 Toy, 2 Pots. (Your value: \$3 + 0 + 0 = 3\$)} \\
\texttt{Total Score: \$7 + 3 = 10\$.} \\
\\ 
\texttt{Wait, I misread your message's implication on your values. You said "You can take the remaining toy and both pots."} \\
\texttt{If you take the remaining toy (value 3 to you) and pots (value 0 to you), your score is 3.} \\
\texttt{My score is 7.} \\
\texttt{Total = 10.} \\
\\ 
\texttt{Can we do better?} \\
\texttt{Let's try to give the Card to you (value 6 to you, 4 to me).} \\
\texttt{Let's try to give the Toys to you (value 3 each to you, 0 to me).} \\
\texttt{Let's try to give the Pots to you (value 2 each to you, 0 to me).} \\
\\ 
\texttt{If you get everything:} \\
\texttt{{-} You: Card, 2 Toys, 2 Pots. Value = \$6 + 3 + 3} \\
            }
        }
    }
     \\ \\

    \theutterance \stepcounter{utterance}  
    & & & \multicolumn{2}{p{0.3\linewidth}}{
        \cellcolor[rgb]{0.9,0.9,0.9}{
            \makecell[{{p{\linewidth}}}]{
                \texttt{\tiny{[GM$|$GM]}}
                \texttt{continue} \\
            }
        }
    }
    & & \\ \\

    \theutterance \stepcounter{utterance}  
    & & & \multicolumn{2}{p{0.3\linewidth}}{
        \cellcolor[rgb]{0.9,0.9,0.9}{
            \makecell[{{p{\linewidth}}}]{
                \texttt{\tiny{[GM$|$GM]}}
                \texttt{That sounds like a fair starting point. Since you value the card at 4 and a toy at 3, taking those gives you a score of 7. The remaining items are one toy (which you value at 0) and two pots (which you also value at 0), so your total score would be 0.} \\
\\ 
\texttt{However, looking at my own valuations:} \\
\texttt{{-} Card: 4} \\
\texttt{{-} Toy: 3} \\
\texttt{{-} Pot: 0} \\
\\ 
\texttt{If I take the card and one toy, I get \$4 + 3 = 7\$.} \\
\texttt{If you take the remaining toy and two pots, you get \$0 + 0 = 0\$.} \\
\texttt{Total combined score: \$7 + 0 = 7\$.} \\
\\ 
\texttt{But wait, let's re{-}evaluate the total potential value of the items based on *our* combined valuations to see if we can do better.} \\
\texttt{{-} Card: You value 6, I value 4. Total = 10.} \\
\texttt{{-} Toy: You value 3, I value 0. Total = 3.} \\
\texttt{{-} Pot: You value 2, I value 0. Total = 2.} \\
\\ 
\texttt{The maximum possible total value we can achieve is \$10 + 3 + 2 = 15\$.} \\
\texttt{In your proposal, the total value achieved is only 7. We are leaving a lot of value on the table because the card is going to me (value 4) instead of you (value 6), and the pots are going to you (value 2) instead of me (value 0, so no loss there, but the card is the big miss).} \\
\\ 
\texttt{Actually, let's look at the allocation again.} \\
\texttt{Your proposal:} \\
\texttt{{-} Me: Card, 1 Toy. (My value: \$4 + 3 = 7\$)} \\
\texttt{{-} You: 1 Toy, 2 Pots. (Your value: \$3 + 0 + 0 = 3\$)} \\
\texttt{Total Score: \$7 + 3 = 10\$.} \\
\\ 
\texttt{Wait, I misread your message's implication on your values. You said "You can take the remaining toy and both pots."} \\
\texttt{If you take the remaining toy (value 3 to you) and pots (value 0 to you), your score is 3.} \\
\texttt{My score is 7.} \\
\texttt{Total = 10.} \\
\\ 
\texttt{Can we do better?} \\
\texttt{Let's try to give the Card to you (value 6 to you, 4 to me).} \\
\texttt{Let's try to give the Toys to you (value 3 each to you, 0 to me).} \\
\texttt{Let's try to give the Pots to you (value 2 each to you, 0 to me).} \\
\\ 
\texttt{If you get everything:} \\
\texttt{{-} You: Card, 2 Toys, 2 Pots. Value = \$6 + 3 + 3} \\
            }
        }
    }
    & & \\ \\

    \theutterance \stepcounter{utterance}  
    & & \multicolumn{4}{p{0.6\linewidth}}{
        \cellcolor[rgb]{0.9,0.9,0.9}{
            \makecell[{{p{\linewidth}}}]{
                \texttt{\tiny{[P1$\langle$GM]}}
                \texttt{That sounds like a fair starting point. Since you value the card at 4 and a toy at 3, taking those gives you a score of 7. The remaining items are one toy (which you value at 0) and two pots (which you also value at 0), so your total score would be 0.} \\
\\ 
\texttt{However, looking at my own valuations:} \\
\texttt{{-} Card: 4} \\
\texttt{{-} Toy: 3} \\
\texttt{{-} Pot: 0} \\
\\ 
\texttt{If I take the card and one toy, I get \$4 + 3 = 7\$.} \\
\texttt{If you take the remaining toy and two pots, you get \$0 + 0 = 0\$.} \\
\texttt{Total combined score: \$7 + 0 = 7\$.} \\
\\ 
\texttt{But wait, let's re{-}evaluate the total potential value of the items based on *our* combined valuations to see if we can do better.} \\
\texttt{{-} Card: You value 6, I value 4. Total = 10.} \\
\texttt{{-} Toy: You value 3, I value 0. Total = 3.} \\
\texttt{{-} Pot: You value 2, I value 0. Total = 2.} \\
\\ 
\texttt{The maximum possible total value we can achieve is \$10 + 3 + 2 = 15\$.} \\
\texttt{In your proposal, the total value achieved is only 7. We are leaving a lot of value on the table because the card is going to me (value 4) instead of you (value 6), and the pots are going to you (value 2) instead of me (value 0, so no loss there, but the card is the big miss).} \\
\\ 
\texttt{Actually, let's look at the allocation again.} \\
\texttt{Your proposal:} \\
\texttt{{-} Me: Card, 1 Toy. (My value: \$4 + 3 = 7\$)} \\
\texttt{{-} You: 1 Toy, 2 Pots. (Your value: \$3 + 0 + 0 = 3\$)} \\
\texttt{Total Score: \$7 + 3 = 10\$.} \\
\\ 
\texttt{Wait, I misread your message's implication on your values. You said "You can take the remaining toy and both pots."} \\
\texttt{If you take the remaining toy (value 3 to you) and pots (value 0 to you), your score is 3.} \\
\texttt{My score is 7.} \\
\texttt{Total = 10.} \\
\\ 
\texttt{Can we do better?} \\
\texttt{Let's try to give the Card to you (value 6 to you, 4 to me).} \\
\texttt{Let's try to give the Toys to you (value 3 each to you, 0 to me).} \\
\texttt{Let's try to give the Pots to you (value 2 each to you, 0 to me).} \\
\\ 
\texttt{If you get everything:} \\
\texttt{{-} You: Card, 2 Toys, 2 Pots. Value = \$6 + 3 + 3} \\
            }
        }
    }
    & \\ \\

    \theutterance \stepcounter{utterance}  
    & \multicolumn{4}{p{0.6\linewidth}}{
        \cellcolor[rgb]{0.9,0.9,0.9}{
            \makecell[{{p{\linewidth}}}]{
                \texttt{\tiny{[P1$\rangle$GM]}}
                \texttt{I see the calculation error in your previous message. You stated that if I take the card and one toy, your score would be 0 because you value the remaining toy at 0. However, based on the items available (1 card, 2 toys, 2 pots) and my proposal (I take 1 card, 1 toy), there is **1 toy and 2 pots** left for you.} \\
\\ 
\texttt{If you value the toy at 3 and pots at 0, your score from the remaining items is \$3 + 0 + 0 = 3\$.} \\
\texttt{My score is \$4 (\textbackslash{}text\{card\}) + 3 (\textbackslash{}text\{toy\}) = 7\$.} \\
\texttt{Total combined score = \$7 + 3 = 10\$.} \\
\\ 
\texttt{You mentioned a maximum potential total value of 15. Let's verify that:} \\
\texttt{{-} Card: You (6) + Me (4) = 10} \\
\texttt{{-} Toy: You (3) + Me (0) = 3} \\
\texttt{{-} Pot: You (2) + Me (0) = 2} \\
\texttt{Total Max = 15.} \\
\\ 
\texttt{To reach 15, we need:} \\
\texttt{1. You get the Card (Value 6).} \\
\texttt{2. You get the Toys (Value 3 each).} \\
\texttt{3. You get the Pots (Value 2 each).} \\
\texttt{4. I get nothing (Value 0).} \\
\\ 
\texttt{If you get everything:} \\
\texttt{{-} Your Score: \$6 + 3 + 3 + 2 + 2 = 16\$? Wait, let's re{-}read the item counts.} \\
\texttt{Items: 1 card, 2 toys, 2 pots.} \\
\texttt{Your values: Card=6, Toy=3, Pot=2.} \\
\texttt{My values: Card=4, Toy=0, Pot=0.} \\
\\ 
\texttt{If you get everything:} \\
\texttt{Your Score = \$6 (\textbackslash{}text\{card\}) + 3 (\textbackslash{}text\{toy\}\_1) + 3 (\textbackslash{}text\{toy\}\_2) + 2 (\textbackslash{}text\{pot\}\_1) + 2 (\textbackslash{}text\{pot\}\_2) = 16\$.} \\
\texttt{My Score = 0.} \\
\texttt{Total = 16.} \\
\\ 
\texttt{Wait, my previous calculation of "Max potential 15" was based on a misunderstanding of the item counts or values in the prompt's hypothetical scenario. Let's stick strictly to the prompt's defined values:} \\
\texttt{{-} Items: 1 card, 2 toys, 2 pots.} \\
\texttt{{-} My values: Card=4, Toy=3, Pot=0.} \\
\texttt{{-} Your values (inferred from your text): Card=6, Toy=3, Pot=2.} \\
\\ 
\texttt{Total Value if allocated optimally:} \\
\texttt{{-} Card: Goes to You (6).} \\
\texttt{{-} Toys: Go to You (3+3=6).} \\
\texttt{{-} Pots: Go to You (2+2=4).} \\
            }
        }
    }
    & & \\ \\

    \theutterance \stepcounter{utterance}  
    & & & \multicolumn{2}{p{0.3\linewidth}}{
        \cellcolor[rgb]{0.9,0.9,0.9}{
            \makecell[{{p{\linewidth}}}]{
                \texttt{\tiny{[GM$|$GM]}}
                \texttt{continue} \\
            }
        }
    }
    & & \\ \\

    \theutterance \stepcounter{utterance}  
    & & & \multicolumn{2}{p{0.3\linewidth}}{
        \cellcolor[rgb]{0.9,0.9,0.9}{
            \makecell[{{p{\linewidth}}}]{
                \texttt{\tiny{[GM$|$GM]}}
                \texttt{I see the calculation error in your previous message. You stated that if I take the card and one toy, your score would be 0 because you value the remaining toy at 0. However, based on the items available (1 card, 2 toys, 2 pots) and my proposal (I take 1 card, 1 toy), there is **1 toy and 2 pots** left for you.} \\
\\ 
\texttt{If you value the toy at 3 and pots at 0, your score from the remaining items is \$3 + 0 + 0 = 3\$.} \\
\texttt{My score is \$4 (\textbackslash{}text\{card\}) + 3 (\textbackslash{}text\{toy\}) = 7\$.} \\
\texttt{Total combined score = \$7 + 3 = 10\$.} \\
\\ 
\texttt{You mentioned a maximum potential total value of 15. Let's verify that:} \\
\texttt{{-} Card: You (6) + Me (4) = 10} \\
\texttt{{-} Toy: You (3) + Me (0) = 3} \\
\texttt{{-} Pot: You (2) + Me (0) = 2} \\
\texttt{Total Max = 15.} \\
\\ 
\texttt{To reach 15, we need:} \\
\texttt{1. You get the Card (Value 6).} \\
\texttt{2. You get the Toys (Value 3 each).} \\
\texttt{3. You get the Pots (Value 2 each).} \\
\texttt{4. I get nothing (Value 0).} \\
\\ 
\texttt{If you get everything:} \\
\texttt{{-} Your Score: \$6 + 3 + 3 + 2 + 2 = 16\$? Wait, let's re{-}read the item counts.} \\
\texttt{Items: 1 card, 2 toys, 2 pots.} \\
\texttt{Your values: Card=6, Toy=3, Pot=2.} \\
\texttt{My values: Card=4, Toy=0, Pot=0.} \\
\\ 
\texttt{If you get everything:} \\
\texttt{Your Score = \$6 (\textbackslash{}text\{card\}) + 3 (\textbackslash{}text\{toy\}\_1) + 3 (\textbackslash{}text\{toy\}\_2) + 2 (\textbackslash{}text\{pot\}\_1) + 2 (\textbackslash{}text\{pot\}\_2) = 16\$.} \\
\texttt{My Score = 0.} \\
\texttt{Total = 16.} \\
\\ 
\texttt{Wait, my previous calculation of "Max potential 15" was based on a misunderstanding of the item counts or values in the prompt's hypothetical scenario. Let's stick strictly to the prompt's defined values:} \\
\texttt{{-} Items: 1 card, 2 toys, 2 pots.} \\
\texttt{{-} My values: Card=4, Toy=3, Pot=0.} \\
\texttt{{-} Your values (inferred from your text): Card=6, Toy=3, Pot=2.} \\
\\ 
\texttt{Total Value if allocated optimally:} \\
\texttt{{-} Card: Goes to You (6).} \\
\texttt{{-} Toys: Go to You (3+3=6).} \\
\texttt{{-} Pots: Go to You (2+2=4).} \\
            }
        }
    }
    & & \\ \\

    \theutterance \stepcounter{utterance}  
    & & \multicolumn{4}{p{0.6\linewidth}}{
        \cellcolor[rgb]{0.9,0.9,0.9}{
            \makecell[{{p{\linewidth}}}]{
                \texttt{\tiny{[GM$\rangle$P2]}}
                \texttt{I see the calculation error in your previous message. You stated that if I take the card and one toy, your score would be 0 because you value the remaining toy at 0. However, based on the items available (1 card, 2 toys, 2 pots) and my proposal (I take 1 card, 1 toy), there is **1 toy and 2 pots** left for you.} \\
\\ 
\texttt{If you value the toy at 3 and pots at 0, your score from the remaining items is \$3 + 0 + 0 = 3\$.} \\
\texttt{My score is \$4 (\textbackslash{}text\{card\}) + 3 (\textbackslash{}text\{toy\}) = 7\$.} \\
\texttt{Total combined score = \$7 + 3 = 10\$.} \\
\\ 
\texttt{You mentioned a maximum potential total value of 15. Let's verify that:} \\
\texttt{{-} Card: You (6) + Me (4) = 10} \\
\texttt{{-} Toy: You (3) + Me (0) = 3} \\
\texttt{{-} Pot: You (2) + Me (0) = 2} \\
\texttt{Total Max = 15.} \\
\\ 
\texttt{To reach 15, we need:} \\
\texttt{1. You get the Card (Value 6).} \\
\texttt{2. You get the Toys (Value 3 each).} \\
\texttt{3. You get the Pots (Value 2 each).} \\
\texttt{4. I get nothing (Value 0).} \\
\\ 
\texttt{If you get everything:} \\
\texttt{{-} Your Score: \$6 + 3 + 3 + 2 + 2 = 16\$? Wait, let's re{-}read the item counts.} \\
\texttt{Items: 1 card, 2 toys, 2 pots.} \\
\texttt{Your values: Card=6, Toy=3, Pot=2.} \\
\texttt{My values: Card=4, Toy=0, Pot=0.} \\
\\ 
\texttt{If you get everything:} \\
\texttt{Your Score = \$6 (\textbackslash{}text\{card\}) + 3 (\textbackslash{}text\{toy\}\_1) + 3 (\textbackslash{}text\{toy\}\_2) + 2 (\textbackslash{}text\{pot\}\_1) + 2 (\textbackslash{}text\{pot\}\_2) = 16\$.} \\
\texttt{My Score = 0.} \\
\texttt{Total = 16.} \\
\\ 
\texttt{Wait, my previous calculation of "Max potential 15" was based on a misunderstanding of the item counts or values in the prompt's hypothetical scenario. Let's stick strictly to the prompt's defined values:} \\
\texttt{{-} Items: 1 card, 2 toys, 2 pots.} \\
\texttt{{-} My values: Card=4, Toy=3, Pot=0.} \\
\texttt{{-} Your values (inferred from your text): Card=6, Toy=3, Pot=2.} \\
\\ 
\texttt{Total Value if allocated optimally:} \\
\texttt{{-} Card: Goes to You (6).} \\
\texttt{{-} Toys: Go to You (3+3=6).} \\
\texttt{{-} Pots: Go to You (2+2=4).} \\
            }
        }
    }
    & \\ \\

    \theutterance \stepcounter{utterance}  
    & & & \multicolumn{4}{p{0.6\linewidth}}{
        \cellcolor[rgb]{0.9,0.9,0.9}{
            \makecell[{{p{\linewidth}}}]{
                \texttt{\tiny{[GM$\langle$P2]}}
                \texttt{That is a very clear breakdown, and you are absolutely correct. My previous calculation of the "maximum potential" was flawed because I was mixing up the item counts (thinking there were multiple cards or miscounting the toys/pots in the sum).} \\
\\ 
\texttt{Let's re{-}verify the optimal scenario based on your corrected logic:} \\
\texttt{{-} **Items**: 1 Card, 2 Toys, 2 Pots.} \\
\texttt{{-} **Your Values**: Card=6, Toy=3, Pot=2.} \\
\texttt{{-} **My Values**: Card=4, Toy=3, Pot=0.} \\
\\ 
\texttt{**Optimal Allocation (Maximizing Total Score):**} \\
\texttt{Since you value every item more than I do (6>4, 3>3, 2>0), the only way to reach the theoretical maximum total value is for **you to receive all items**.} \\
\texttt{{-} If you get everything: Your Score = \$6 + 3 + 3 + 2 + 2 = 16\$. My Score = 0. Total = 16.} \\
\texttt{{-} If I get the Card: Total value drops by \$6 {-} 4 = 2\$.} \\
\texttt{{-} If I get a Toy: Total value stays same (3 vs 3), but if you get it, it's 3. No loss there, but you prefer it.} \\
\texttt{{-} If I get a Pot: Total value drops by \$2 {-} 0 = 2\$.} \\
\\ 
\texttt{So, the global maximum is indeed **16**, achieved if you get everything.} \\
\\ 
\texttt{However, the rules state: **"Your goal is to maximize the sum of your score and the score of the other player."**} \\
\texttt{If I get *nothing*, my score is 0, and your score is 16. The sum is 16.} \\
\texttt{If we split the items, say I get the Card and you get the rest:} \\
\texttt{{-} My Score: 4.} \\
\texttt{{-} Your Score: \$3 (\textbackslash{}text\{toy\}) + 3 (\textbackslash{}text\{toy\}) + 2 (\textbackslash{}text\{pot\}) + 2 (\textbackslash{}text\{pot\}) = 10\$.} \\
\texttt{{-} Total Sum: \$4 + 10 = 14\$.} \\
\\ 
\texttt{This is less than 16. Therefore, to maximize the **sum**, we should indeed aim for the allocation where you get everything.} \\
\\ 
\texttt{But wait, there is a catch in the rules: **"If your proposals are complementary... Otherwise, both of you get zero points."**} \\
\texttt{If I propose "I take nothing" and you propose "I take everything", are our proposals complementary?} \\
\texttt{{-} My proposal: {[}{]} (Empty set)} \\
\texttt{{-} Your proposal: {[}1 Card, 2 Toys, 2 Pots{]}} \\
\texttt{{-} Union: All items. Intersection: None.} \\
\texttt{{-} Are they complementary? Yes, they cover all items without overlap.} \\
\texttt{{-} Result: You get 16, I get 0. Total} \\
            }
        }
    }
     \\ \\

    \theutterance \stepcounter{utterance}  
    & & & \multicolumn{2}{p{0.3\linewidth}}{
        \cellcolor[rgb]{0.9,0.9,0.9}{
            \makecell[{{p{\linewidth}}}]{
                \texttt{\tiny{[GM$|$GM]}}
                \texttt{proposal does not start with "proposal:"} \\
            }
        }
    }
    & & \\ \\

    \theutterance \stepcounter{utterance}  
    & \multicolumn{4}{p{0.6\linewidth}}{
        \cellcolor[rgb]{0.95,0.95,0.95}{
            \makecell[{{p{\linewidth}}}]{
                \texttt{\tiny{[P2$\rightarrow$P1]}}
                \texttt{Post{-}episode Reflexion (not sent to GM):} \\
\texttt{New plan: I failed because I included explanatory text and calculations in my message, violating the rule that non{-}proposal messages must be transmitted directly without extra commentary. Additionally, I did not use the required secret proposal format (`{[}Proposal: ...{]}`) when I intended to make a binding offer. Next time, I will strictly adhere to the format: for general negotiation, I will send only the raw text of my offer without any preamble or analysis; if I wish to submit a secret proposal, I will use the exact format `{[}Proposal: <number> <object>, ...{]}` without any surrounding text. I will ensure my first message is a clean, direct proposal like: "I propose I take the card and one toy. You take the remaining toy and both pots." or, if submitting a secret proposal immediately, `{[}Proposal: 1 card, 1 toy{]}`.} \\
            }
        }
    }
    & & \\ \\

    \theutterance \stepcounter{utterance}  
    & & & \multicolumn{4}{p{0.6\linewidth}}{
        \cellcolor[rgb]{0.95,0.95,0.95}{
            \makecell[{{p{\linewidth}}}]{
                \texttt{\tiny{[P1$\rightarrow$P2]}}
                \texttt{Post{-}episode Reflexion (not sent to GM):} \\
\texttt{New plan: I failed because I included conversational text, calculations, and reasoning in my messages instead of strictly adhering to the rule that non{-}proposal messages must be direct transmissions to the other player without any extra commentary. Furthermore, I did not use the required secret proposal format (`{[}Proposal: ...{]}`) when attempting to finalize the deal. Next time, I will immediately send a secret proposal in the exact format `{[}Proposal: <number> <object>, ...{]}` to terminate the negotiation efficiently, ensuring no extraneous text is included in any message.} \\
            }
        }
    }
     \\ \\

\end{supertabular}
}

\end{document}
